\chapter{Introduction}

\section{Course Purpose and Learning Objectives}

This introductory chapter explains how the \textbf{Applied Physics in Industrial Design (APID)} project will run, what your team is expected to deliver, and how the hydroponic theme connects design work to applied physics. In this session, you will review the course workflow, assessment structure, teamwork expectations, and support resources before moving into later technical sessions.

By the end of this introduction, you should be able to:
\begin{enumerate}
    \item explain the purpose of the APID project and its project-based learning approach;
    \item describe why hydroponic agriculture is the shared project theme for the semester;
    \item identify the main deliverables, submission routes, and assessment formats used in the module; and
    \item organize your team so that weekly tasks, records, and communication remain manageable throughout the semester.
\end{enumerate}

Table~\ref{tab:session0-key-terms} defines the main terms used throughout this chapter.

\begin{table}[htbp]
\centering
\caption{Key terms introduced in the APID orientation}
\label{tab:session0-key-terms}
\renewcommand{\arraystretch}{1.2}
\begin{tabularx}{\textwidth}{|>{\raggedright\arraybackslash}p{0.2\textwidth}|X|}
\hline
\textbf{Term} & \textbf{Definition} \\
\hline
APID & Applied Physics in Industrial Design, the course that combines product development practice with applied physics reasoning. \\
\hline
Hydroponic agriculture & A method of growing plants without soil by controlling water, nutrients, support structures, and environmental conditions. \\
\hline
LLM & Large Language Model, a generative artificial intelligence tool that can assist with drafting, summarizing, or brainstorming but must be checked critically. \\
\hline
CLO & Course Learning Outcome, the formal learning outcome used to structure assessment tasks. \\
\hline
ITEL & Institutional Teaching and E-Learning platform, the university system used for announcements, submission links, and updated deadlines. \\
\hline
PD & Product development, the structured process of turning an opportunity into a product concept, evaluation plan, and communication package. \\
\hline
\end{tabularx}
\end{table}

\section{Project Theme: Hydroponic Agriculture}

This year marks the third time APID has been offered. Previously, the module used a more open-ended project theme, which allowed for a wide variety of product ideas. However, the open theme also made it difficult for students to find a common context for applied physics concepts, and it led to some confusion about how to connect design work to technical reasoning. To address these issues, the project theme is now focused on hydroponic agriculture, which provides a shared context for all teams while still allowing for creativity and diversity in product ideas. 

Therefore, your team will propose a product related to hydroponic agriculture. The product may operate directly inside a hydroponic system or support hydroponic users, growers, or maintenance workflows.

Just to reiterate, hydroponics is a suitable project theme because it requires you to consider several applied physics topics simultaneously, including but not limited to:
\begin{itemize}
    \item fluid transport and pressure losses in nutrient delivery systems;
    \item heat transfer and environmental control for plant growth;
    \item sensing and instrumentation for pH, electrical conductivity, temperature, or flow monitoring; and
    \item materials selection, structural stability, and manufacturability for safe, durable products.
\end{itemize}

The theme therefore keeps the project grounded in industrial design while still demanding technical reasoning. Your task is not merely to propose a visually interesting idea. You are expected to justify why the concept should work, who it serves, and how it can be developed systematically.

\section{Team Formation, Roles, and Working Method}

Teams are assigned by the module convenor, usually with four or five students per group. Random grouping is used to create diversity of experience and to simulate professional collaboration with unfamiliar teammates.

Once the group list is published on ITEL, review it immediately. If your name is missing, if you are an exchange student without an assigned team, or if a listed member is no longer attending the module, contact the module convenor as soon as possible.

Each group should appoint a leader. To ensure fairness, the group leader role should be rotated every one or two weeks. The leader will play a management role in addition to their normal work as a group member. This includes ensuring that everyone is involved, work is divided fairly, tasks are completed on time, and project milestones are met. The leader does not have to make all the decisions but has a responsibility to ensure the group works effectively as a whole, keeping everyone aligned with the project goals and deadlines. Once assigned to a project team, you are expected to remain in that team for the entire term.

\section{Weekly Workflow and Meeting Expectations}

All sessions consist of a \textbf{pre-class task} and an \textbf{in-class task}. 

Pre-class activities refer to tasks that must be conducted by each member prior to attending the tutorial session. These activities are designed to help you prepare for tutorial discussions and to make progress on the project. They may include reading, research, calculations, or drafting outputs that will be discussed in class. The tutorial discussion is an opportunity to share your work, receive feedback from peers and the lecturer, and identify areas for improvement. 

The \textbf{in-class task} is the activity conducted during the tutorial session. During these sessions, you will compare ideas, receive feedback, and identify what must be improved. After class, your team should refine the material so that it can be used in later deliverables.

It is also possible for groups to conduct meetings outside of class time. These meetings are normally held in a mutually convenient location on campus and have no prescribed time limit. The content of these meetings is determined by the group; they are usually informal, but some groups may find it useful to maintain practices such as appointing a chairperson or taking minutes.

\section{Deliverables and Assessment Overview}

Table~\ref{tab:session0-deliverables} summarizes the major deliverables introduced in this chapter. Exact dates, upload links, and any weighting updates must always be confirmed through ITEL.

\begin{table}[htbp]
\centering
\caption{Major APID deliverables introduced in Session 0}
\label{tab:session0-deliverables}
\renewcommand{\arraystretch}{1.2}
\begin{tabularx}{\textwidth}{|>{\raggedright\arraybackslash}p{0.2\textwidth}|>{\raggedright\arraybackslash}p{0.18\textwidth}|>{\raggedright\arraybackslash}p{0.2\textwidth}|X|c|}
\hline
\textbf{Deliverable} & \textbf{Typical timing} & \textbf{Submission format} & \textbf{Purpose} & \textbf{Weight (\%)}\\
\hline
Weekly tutorial exercises & Throughout the semester & PDF or specified format via ITEL & Documents progress in needs analysis, specifications, and concept generation for the final report and presentation. & 0 \\
\hline
Presentation Day & End-of-semester period & Slides and optional demonstration & Communicates the problem, solution, and reflective learning outcomes. & 20 \\
\hline
Final Report & End of semester & Single group PDF via ITEL & Integrates the full product development process into a self-contained technical report. & 40 \\
\hline
Final Exam & Exam period & 3-hour exam & Tests understanding of Concept Selection, Testing, Industrial Design, and IP. & 40 \\
\hline
\end{tabularx}
\end{table}

\section{Resources, Support, and Project Constraints}

This section highlights the practical constraints affecting the project. The module does not provide a dedicated project budget, so teams should choose scopes and materials realistically. If your concept requires prototype expenses, those costs must be managed by the team unless a different arrangement is approved by the lecturer.

\section*{Intellectual Property Rights}

In this course, any inventions made by your team will usually belong to the team members. If you decide to apply for a patent, you can either do it independently (bearing all related costs) or choose to patent it with the help of Universiti Malaysia Sabah. In the latter case, the university will cover expenses and necessary steps. It is important to discuss and agree with your team members early on about how to divide any potential earnings. Additionally, any ideas shared on digital whiteboards or similar tools should be preserved as recorded evidence of the invention's development.
